\documentclass[11pt]{article}
\usepackage{amsmath, amssymb, amsthm}

\title{A6 Unified Kernel Geometry}
\author{Borealis S. Hedling}
\date{October 2026}

\begin{document}
\maketitle

\section*{Abstract}
This document defines the unified kernel for A6 dual-algebra systems. The kernel
acts as the geometric firewall binding spectral, modal, holonomy, collapse, and
stability operators into a single invariant-preserving structure. It provides
the regular-value constraints, tangent-space structure, intersection geometry,
and provenance lift required for safe operator composition and runtime
diagnostics.

\section{Kernel Geometry}
Let $X$ be the base geometric manifold defined in A6GeometryDefinitions. The
unified kernel $K \subset X$ is defined as the intersection of all operator
kernels:
\[
K = K_{\text{spec}} \cap K_{\text{modal}} \cap K_{\text{holo}} \cap
    K_{\text{collapse}} \cap K_{\text{stab}}.
\]

Each kernel is a regular-value submanifold:
\[
K_{\text{op}} = f_{\text{op}}^{-1}(0),
\]
where $f_{\text{op}}$ is the operator-specific constraint function.

\section{Regular-Value Conditions}
For each operator constraint $f_{\text{op}}$, the differential
$Df_{\text{op}}(p)$ is surjective for all $p \in K_{\text{op}}$. This ensures
that each kernel is a smooth submanifold and that the unified kernel is a
transversal intersection.

\section{Tangent-Space Structure}
The tangent space of the unified kernel is given by:
\[
T_p K = \bigcap_{\text{op}} \ker Df_{\text{op}}(p).
\]
This intersection is non-degenerate due to the regular-value conditions and
transversality of the operator constraints.

\section{Intersection Geometry}
The unified kernel inherits geometric structure from each operator kernel. In
particular:
\begin{itemize}
  \item spectral continuity restricts admissible tangent directions,
  \item modal adjacency restricts allowable transitions,
  \item holonomy drift restricts curvature accumulation,
  \item collapse surfaces restrict admissible neighborhoods,
  \item stability envelopes restrict allowable perturbations.
\end{itemize}

\section{Dual-Algebra Firewall}
The unified kernel acts as a firewall preventing operator divergence. For any
operator $O$ and point $p \in K$:
\[
O(p) \in K.
\]
This ensures that operator composition preserves kernel membership and prevents
catastrophic geometry.

\section{Provenance Lift}
Each operator has a provenance predicate $P_{\text{op}}$. The unified kernel
provenance is defined as:
\[
P_K(p) = \bigwedge_{\text{op}} P_{\text{op}}(p).
\]
This ensures that provenance is preserved across operator composition.

\section{Collapse-Surface Interaction}
Collapse surfaces $\Sigma$ intersect the unified kernel transversally:
\[
T_p K + T_p \Sigma = T_p X.
\]
This ensures that collapse detection is well-defined and that collapse surfaces
do not degenerate inside the kernel.

\section{Stability-Envelope Interaction}
The stability envelope $E$ restricts admissible perturbations:
\[
E(p) \leq E(q) \quad \text{for all } p, q \in K.
\]
This ensures that kernel perturbations remain bounded.

\section{Operator-Kernel Binding}
For each operator $O$:
\[
O : K \to K.
\]
This binding ensures that operator semantics preserve kernel membership and that
runtime diagnostics can safely monitor operator behavior.

\section*{Conclusion}
The unified kernel provides the geometric foundation required for safe operator
composition, invariant preservation, collapse detection, stability envelopes,
and runtime diagnostics. It completes the A6 physics layer and enables
construction of the diagnostic and safety layers required for A7 altitude.

\end{document}
